\section{External Benchmark Evaluation: MEHDIE}

This section extends the evaluation with experiments on the MEHDIE benchmark.
The MEHDIE datasets were not a design target of the proposed model; results are therefore
interpreted as a stress-test of cross-corpus robustness rather than as a primary optimisation goal.

\subsection{Benchmark Scope and Interpretation}

The proposed model is designed for large-scale gazetteer reconciliation, with GeoNames-scale
candidate sets and extreme class imbalance. In contrast, MEHDIE focuses on medieval manuscript
alignment with carefully curated phonetic and orthographic assumptions. Direct performance parity
with MEHDIE systems is therefore neither expected nor required for validation of the model’s core aims.

\subsection{Results Overview}

Table~\ref{tab:mehdie_f5_best} reports the best F$_5$ score achieved by each method on the five MEHDIE
testsets. The F$_5$ metric weights recall five times more heavily than precision, reflecting user
preference for exhaustive candidate discovery.

\begin{table}[h]
\centering
\begin{tabular}{lccc}
\hline
Testset & Levenshtein & Jaro--Winkler & Our Model \\
\hline
TS7 & 0.476 & 0.410 & 0.371 \\
TS8 & 0.750 & 0.653 & 0.738 \\
TS9 & 0.768 & 0.731 & 0.676 \\
TS10 & 0.428 & 0.411 & 0.437 \\
TS11 & 0.832 & 0.784 & 0.857 \\
\hline
Average & 0.651 & 0.598 & 0.616 \\
\hline
\end{tabular}
\caption{Best F$_5$ scores per MEHDIE testset.}
\label{tab:mehdie_f5_best}
\end{table}

\subsection{Comparison with Published MEHDIE Results}

Table~\ref{tab:mehdie_comparison} compares the proposed model against the best published MEHDIE
orthographic or phonetic result for each testset.

\begin{table}[h]
\centering
\begin{tabular}{lccc}
\hline
Testset & MEHDIE Best & Our Model & $\Delta$ \\
\hline
TS7 & 0.77 & 0.371 & -0.399 \\
TS8 & 0.68 & 0.738 & +0.058 \\
TS9 & 0.92 & 0.676 & -0.244 \\
TS10 & 0.77 & 0.437 & -0.333 \\
TS11 & 0.88 & 0.857 & -0.023 \\
\hline
\end{tabular}
\caption{Comparison with MEHDIE published F$_5$ scores.}
\label{tab:mehdie_comparison}
\end{table}

\subsection{Discussion}

The results demonstrate that the proposed model generalises beyond its training distribution,
matching or exceeding classical string baselines on several MEHDIE testsets and approaching
specialised MEHDIE phonetic performance in selected cases (notably TS8 and TS11).
Where performance is substantially lower (TS7, TS9), the datasets strongly favour handcrafted
phonetic assumptions specific to the MEHDIE framework.

These findings support the conclusion that the learned phonetic embedding space is robust and
transferable, while reinforcing the distinction between candidate generation and final decision-making.

\subsection{Future Work}

Several extensions are planned:
\begin{itemize}
  \item Primary evaluation will continue to foreground GeoNames-scale experiments, which directly
        reflect the model’s intended use.
  \item OpenStreetMap-based evaluation is a promising next step due to noisier labels and richer aliasing.
  \item MEHDIE datasets could be incorporated as additional fine-tuning or curriculum training data.
  \item A ``brutal'' ablation study is planned to isolate the contributions of phonetic normalisation,
        sequence sensitivity, and metric learning.
\end{itemize}